\chapter{Weighted Residual Methods}
\label{ch:weighted_residuals}

\section{Motivation}

The finite element method can be introduced from several points of view:
direct stiffness, energy methods, virtual work, or weighted residual methods.
The weighted residual viewpoint is useful because it shows how an approximate
solution is obtained when the differential equation cannot be satisfied exactly
at every point of the domain.

Consider a trial solution $\tilde{u}(x)$ that satisfies the essential boundary
conditions but is not the exact solution. After substitution into the governing
differential equation, a nonzero residual remains:
\begin{equation}
    R(x) = \mathcal{L}\tilde{u}(x) - f(x)
    \label{eq:residual_definition}
\end{equation}
where $\mathcal{L}$ is the differential operator and $f(x)$ is the loading term.
Weighted residual methods determine the unknown coefficients in $\tilde{u}$ by
requiring the residual to vanish in an averaged or weighted sense:
\begin{equation}
    \int_{\Omega} w_i(x) R(x) \dd \Omega = 0,
    \qquad i = 1,2,\ldots,n
    \label{eq:weighted_residual_statement}
\end{equation}
where $w_i(x)$ are weighting functions. Different choices of $w_i$ lead to
different methods.

\section{Axially Loaded Bar}

Consider the prismatic member in \cref{fig:wr_bar}, fixed at $x=0$ and free at
$x=L$. Since the load acts only along the axis of the member, the structural
model is a bar or rod rather than a bending beam. The bar is subjected to a
distributed axial load
\begin{equation}
    p(x) = p_0 \left(\frac{x}{L}\right)^2
    \label{eq:wr_distributed_load}
\end{equation}
with constant $E$ and $A$.

\begin{figure}[H]
\centering
\begin{tikzpicture}[>=Stealth, x=1cm, y=1cm]
  \tikzset{
    load/.style={->,thick,blue!70!black},
    dim/.style={<->,thin},
    tick/.style={thin},
  }
  %% Fixed wall
  \fill[pattern=north east lines, pattern color=gray!70] (-0.32,-0.42) rectangle (0,0.42);
  \draw[thick] (0,-0.42) -- (0,0.42);
  %% Bar body
  \fill[gray!18] (0,-0.16) rectangle (6.2,0.16);
  \draw[thick] (0,-0.16) rectangle (6.2,0.16);
  %% Distributed axial load arrows, increasing to the right
  \foreach \x/\len in {0.65/0.30,1.30/0.38,1.95/0.46,2.60/0.54,3.25/0.62,3.90/0.70,4.55/0.78,5.20/0.86}{
    \draw[load] (\x,0.55) -- ++(\len,0);
  }
  \node[blue!70!black,above] at (3.25,0.72) {$p(x)=p_0(x/L)^2$};
  %% Free end label
  \node[anchor=west] at (6.42,0.18) {free end};
  %% Dimension and axis
  \draw[tick] (0,-0.52) -- (0,-0.82);
  \draw[tick] (6.2,-0.52) -- (6.2,-0.82);
  \draw[dim] (0,-0.68) -- (6.2,-0.68)
    node[midway,fill=white,inner sep=1pt] {$L$};
  \node[below] at (0,-0.82) {\small $x=0$};
  \node[below] at (6.2,-0.82) {\small $x=L$};
  \draw[->,thin] (6.55,-0.68) -- (7.20,-0.68) node[right] {$x$};
\end{tikzpicture}
\caption{Fixed-free bar subjected to a nonuniform axial load.}
\label{fig:wr_bar}
\end{figure}

\subsection{Equilibrium of a Differential Element}

To derive the governing equation, isolate a differential element of length
$\dd x$ as shown in \cref{fig:wr_differential_element}. Let $N(x)$ denote the
internal axial force resultant. The sign convention is that $p(x)$ is positive
in the positive $x$-direction and $N$ is positive in tension.

\begin{figure}[H]
\centering
\begin{tikzpicture}[>=Stealth, x=1cm, y=1cm]
  \tikzset{
    load/.style={->,thick,blue!70!black},
    force/.style={->,very thick,red!70!black},
    dim/.style={<->,thin},
    cut/.style={dashed,gray!65},
    tick/.style={thin},
  }
  %% Differential element
  \fill[gray!18] (1.55,-0.25) rectangle (5.15,0.25);
  \draw[thick] (1.55,-0.25) rectangle (5.15,0.25);
  %% Cut faces
  \draw[cut] (1.55,-0.52) -- (1.55,0.64);
  \draw[cut] (5.15,-0.52) -- (5.15,0.64);
  %% Internal axial forces
  \draw[force] (1.55,0) -- (0.25,0)
    node[midway,above=2pt] {$N(x)$};
  \draw[force] (5.15,0) -- (6.65,0)
    node[midway,above=2pt] {$N(x+\dd x)$};
  %% Equivalent resultant of the distributed axial load over dx
  \draw[load,very thick] (2.75,0.68) -- (4.10,0.68);
  \node[blue!70!black,above] at (3.43,0.86) {$p(x)\dd x$};
  %% Coordinates and dimension
  \draw[tick] (1.55,-0.58) -- (1.55,-0.88);
  \draw[tick] (5.15,-0.58) -- (5.15,-0.88);
  \draw[dim] (1.55,-0.73) -- (5.15,-0.73)
    node[midway,fill=white,inner sep=1pt] {$\dd x$};
  \node[below] at (1.55,-0.88) {\small $x$};
  \node[below] at (5.15,-0.88) {\small $x+\dd x$};
  \draw[->,thin] (5.55,-0.73) -- (6.25,-0.73) node[right] {$x$};
\end{tikzpicture}
\caption{Free-body diagram of a differential bar element.}
\label{fig:wr_differential_element}
\end{figure}

Force equilibrium in the axial direction gives
\begin{equation}
    -N(x) + N(x+\dd x) + p(x)\dd x = 0 .
    \label{eq:wr_force_equilibrium_discrete}
\end{equation}
Using the Taylor expansion
\[
    N(x+\dd x) = N(x) + \frac{\dd N}{\dd x}\dd x
\]
and neglecting higher-order terms, \cref{eq:wr_force_equilibrium_discrete}
becomes
\begin{equation}
    \frac{\dd N}{\dd x} + p(x) = 0 .
    \label{eq:wr_axial_equilibrium}
\end{equation}

\subsection{Constitutive and Kinematic Relations}

The axial strain is related to the displacement field by
\begin{equation}
    \varepsilon(x) = \frac{\dd u}{\dd x}.
    \label{eq:wr_axial_strain}
\end{equation}
For a linear elastic material,
\begin{equation}
    \sigma(x) = E\varepsilon(x).
    \label{eq:wr_hooke_law}
\end{equation}
The axial force resultant is
\begin{equation}
    N(x) = \sigma(x)A = EA\frac{\dd u}{\dd x}.
    \label{eq:wr_axial_force_displacement}
\end{equation}
Substituting \cref{eq:wr_axial_force_displacement} into
\cref{eq:wr_axial_equilibrium} gives the differential equation of an axially
loaded bar:
\begin{equation}
    \frac{\dd}{\dd x}\!\left(EA\frac{\dd u}{\dd x}\right) + p(x) = 0 .
    \label{eq:wr_general_strong_bar}
\end{equation}
Equivalently,
\begin{equation}
    -\frac{\dd}{\dd x}\!\left(EA\frac{\dd u}{\dd x}\right) = p(x).
\end{equation}

For constant $E$ and $A$, and for the load in
\cref{eq:wr_distributed_load}, the governing equation becomes
\begin{align}
    -EA \frac{\dd^2 u}{\dd x^2} &= p_0\left(\frac{x}{L}\right)^2,
    \qquad 0 < x < L
    \label{eq:wr_strong_bar} \\
    u(0) &= 0, \qquad EA u'(L) = 0 .
    \notag
\end{align}
The first boundary condition is essential because it prescribes displacement.
The second is natural because it prescribes axial force. If an external axial
force $F$ were applied at $x=L$, the natural boundary condition would be
$EAu'(L)=F$.

For this introductory comparison, use the one-term trial function
\begin{equation}
    \tilde{u}(x) = a \left[\left(\frac{x}{L}\right)^2
    - 2\left(\frac{x}{L}\right)\right]
    \label{eq:wr_trial_function}
\end{equation}
where $a$ is the unknown coefficient. This function satisfies both boundary
conditions:
\[
    \tilde{u}(0)=0, \qquad EA\tilde{u}'(L)=0 .
\]
It is not flexible enough to satisfy the differential equation exactly, so a
residual remains.

Let
\[
    \xi = \frac{x}{L}, \qquad 0 \leq \xi \leq 1 .
\]
Substituting \cref{eq:wr_trial_function} into \cref{eq:wr_strong_bar} gives
\begin{equation}
    R(\xi) =
    -EA \frac{\dd^2\tilde{u}}{\dd x^2} - p_0\xi^2
    =
    -\frac{2EA}{L^2}a - p_0\xi^2 .
    \label{eq:wr_bar_residual}
\end{equation}
The three methods below differ only in how they choose $a$ from this residual.

\section{Collocation Method}

The collocation method requires the residual to vanish at selected points in
the domain. Since the approximation contains one unknown coefficient, one
collocation point is needed. Choosing the midpoint $\xi=1/2$ gives
\begin{equation}
    R\left(\frac{1}{2}\right)=0 .
\end{equation}
Using \cref{eq:wr_bar_residual},
\begin{equation}
    -\frac{2EA}{L^2}a - \frac{p_0}{4} = 0
    \qquad\Rightarrow\qquad
    a_{\mathrm{col}} = -\frac{p_0L^2}{8EA}.
    \label{eq:wr_collocation_a}
\end{equation}
The approximate displacement is therefore
\begin{equation}
    \tilde{u}_{\mathrm{col}}(x)
    = -\frac{p_0L^2}{8EA}\left(\xi^2-2\xi\right).
\end{equation}

\section{Subdomain Method}

The subdomain method requires the residual to have zero average over each
subdomain. With one unknown coefficient, the whole bar is taken as one
subdomain:
\begin{equation}
    \int_0^L R(x)\dd x = 0 .
\end{equation}
Equivalently,
\begin{equation}
    L\int_0^1 \left(
    -\frac{2EA}{L^2}a - p_0\xi^2
    \right)\dd \xi = 0 .
\end{equation}
Thus,
\begin{equation}
    -\frac{2EA}{L^2}a - \frac{p_0}{3} = 0
    \qquad\Rightarrow\qquad
    a_{\mathrm{sub}} = -\frac{p_0L^2}{6EA}.
    \label{eq:wr_subdomain_a}
\end{equation}

\section{Galerkin Method}

In the Galerkin method, the weighting function is chosen from the same family
as the trial function. Here,
\begin{equation}
    w(\xi) = \xi^2 - 2\xi .
    \label{eq:wr_galerkin_weight}
\end{equation}
The weighted residual statement is
\begin{equation}
    \int_0^L w(x)R(x)\dd x = 0 .
\end{equation}
Substituting \cref{eq:wr_bar_residual,eq:wr_galerkin_weight} and changing to
$\xi$ gives
\begin{equation}
    L\int_0^1 (\xi^2-2\xi)
    \left(
    -\frac{2EA}{L^2}a - p_0\xi^2
    \right)\dd \xi = 0 .
\end{equation}
The required integrals are
\[
    \int_0^1(\xi^2-2\xi)\dd\xi = -\frac{2}{3},
    \qquad
    \int_0^1(\xi^2-2\xi)\xi^2\dd\xi = -\frac{3}{10}.
\]
Therefore,
\begin{equation}
    \frac{4EA}{3L^2}a + \frac{3p_0}{10} = 0
    \qquad\Rightarrow\qquad
    a_{\mathrm{gal}} = -\frac{9p_0L^2}{40EA}.
    \label{eq:wr_galerkin_a}
\end{equation}

\begin{importantbox}
The Galerkin method is the weighted residual method that leads directly to the
standard finite element formulation when piecewise polynomial trial functions
are used and the governing equation is written in weak form.
\end{importantbox}

\section{Comparison with the Exact Solution}

For this example, the exact solution is obtained by integrating
\cref{eq:wr_strong_bar} twice and applying the boundary conditions:
\begin{equation}
    u(x) = \frac{p_0L^2}{EA}
    \left(\frac{\xi}{3} - \frac{\xi^4}{12}\right).
    \label{eq:wr_exact_solution}
\end{equation}
At the free end,
\begin{equation}
    u(L) = \frac{p_0L^2}{4EA}.
\end{equation}
Since $\tilde{u}(L)=-a$, the three approximations give the tip displacements
listed in \cref{tab:wr_comparison}.

\begin{table}[htbp]
\centering
\caption{Comparison of one-term weighted residual approximations.}
\label{tab:wr_comparison}
\begin{tabular}{@{}lll@{}}
    \toprule
    Method & Coefficient $a$ & Tip displacement $\tilde{u}(L)$ \\
    \midrule
    Collocation &
    $-\dfrac{p_0L^2}{8EA}$ &
    $\dfrac{p_0L^2}{8EA}$ \\
    Subdomain &
    $-\dfrac{p_0L^2}{6EA}$ &
    $\dfrac{p_0L^2}{6EA}$ \\
    Galerkin &
    $-\dfrac{9p_0L^2}{40EA}$ &
    $\dfrac{9p_0L^2}{40EA}$ \\
    Exact &
    -- &
    $\dfrac{p_0L^2}{4EA}$ \\
    \bottomrule
\end{tabular}
\end{table}

The collocation method makes the residual zero at one selected point, but the
residual may still be large elsewhere. The subdomain method balances the total
residual over the domain. The Galerkin method weights the residual with the
same function used to approximate the displacement, which gives the best result
among these one-term approximations for this problem.

\section{MATLAB Implementation and Results}

The following MATLAB script evaluates the coefficients, prints a numerical
tip-displacement comparison, and generates the normalized displacement curves
shown in \cref{fig:wr_matlab_comparison}.

\Cref{fig:wr_matlab_comparison} compares the three weighted residual
approximations with the exact solution. The displacement is normalized as
\(uEA/(p_0L^2)\), so the plot shows the method error independently of the
chosen numerical values of \(E\), \(A\), \(p_0\), and \(L\).

\begin{figure}[H]
\centering
\includegraphics[width=0.82\textwidth]{figures/weighted_residual_comparison.pdf}
\caption{MATLAB comparison of collocation, subdomain, and Galerkin one-term
approximations against the exact solution.}
\label{fig:wr_matlab_comparison}
\end{figure}

\begin{lstlisting}[caption={MATLAB comparison of weighted residual methods},
                   label={lst:wr_matlab_comparison}]
% weighted_residual_comparison.m
% One-term weighted residual approximation for an axially loaded bar

clear; clc; close all;

%% Problem parameters
L  = 1.0;        % bar length [m]
E  = 210e9;      % Young's modulus [Pa]
A  = 1.0e-4;     % cross-sectional area [m^2]
p0 = 1.0e4;      % load intensity scale [N/m]
EA = E*A;        % axial stiffness [N]

%% Method coefficients for u_tilde = a*(xi^2 - 2*xi)
a_col = -p0*L^2/(8*EA);       % collocation at xi = 1/2
a_sub = -p0*L^2/(6*EA);       % one subdomain over [0,L]
a_gal = -9*p0*L^2/(40*EA);    % Galerkin weighting

methodNames = {'Collocation'; 'Subdomain'; 'Galerkin'};
aValues = [a_col; a_sub; a_gal];
tipValues = -aValues;         % xi=1 gives u_tilde(L) = -a

%% Exact tip displacement and percent error
uTipExact = p0*L^2/(4*EA);
percentError = abs(tipValues - uTipExact)/uTipExact * 100;

fprintf('Weighted residual comparison for axial bar\n');
fprintf('L = %.3f m, EA = %.4e N, p0 = %.4e N/m\n\n', L, EA, p0);
fprintf('%-14s %14s %18s %12s\n', ...
        'Method', 'a [m]', 'u_tilde(L) [m]', 'Error [%]');
fprintf('%s\n', repmat('-', 1, 62));

for i = 1:numel(methodNames)
    fprintf('%-14s %14.6e %18.6e %12.3f\n', ...
        methodNames{i}, aValues(i), tipValues(i), percentError(i));
end
fprintf('%-14s %14s %18.6e %12.3f\n', 'Exact', '--', uTipExact, 0);

%% Normalized displacement curves
xi = linspace(0, 1, 400);
phi = xi.^2 - 2*xi;

u_col_norm = -(1/8)  * phi;
u_sub_norm = -(1/6)  * phi;
u_gal_norm = -(9/40) * phi;
u_exact_norm = xi/3 - xi.^4/12;

figure('Name','Weighted Residual Comparison','Color','w');
plot(xi, u_exact_norm, 'k-',  'LineWidth', 1.8); hold on;
plot(xi, u_col_norm,   'b--', 'LineWidth', 1.3);
plot(xi, u_sub_norm,   'r-.', 'LineWidth', 1.3);
plot(xi, u_gal_norm,   'g:',  'LineWidth', 1.9);
grid on; box on;
xlabel('Normalized coordinate, \xi = x/L');
ylabel('Normalized displacement, uEA/(p_0L^2)');
legend('Exact', 'Collocation', 'Subdomain', 'Galerkin', ...
       'Location', 'northwest');
\end{lstlisting}

\section{Connection to FEM}

The finite element method uses the Galerkin idea locally on each element. The
unknown displacement field is approximated by shape functions,
\begin{equation}
    u^e(x) = \vect{N}(x)\vect{u}^e ,
\end{equation}
and the weighting functions are chosen as the same shape functions. After
integration by parts, the weak form for the axial bar becomes
\begin{equation}
    \int_0^L EA\, \frac{\dd w}{\dd x}
    \frac{\dd u}{\dd x}\dd x
    =
    \int_0^L w p \dd x + w(L)F .
    \label{eq:wr_weak_form_connection}
\end{equation}
This is the starting point for the bar element stiffness matrix derived in
\cref{ch:1d}. In other words, the familiar finite element equations are not
separate from weighted residual methods; they are the Galerkin weighted residual
method applied element by element.
